\chapter{Аудит существующих кратностей и аффинный резонансный сток}

Минимальный невырожденный холодильник не прошёл, поскольку поперечный
дублет семейного триплета требует двух ортогональных принимающих мод на
одном энергетическом разрыве. Прежде чем добавлять второй канал, настоящий
проверка рассматривает все крупные размерности проекта: $15$, $18$, $20$, $35$,
KO6-удвоение, Real-пару и аффинное меню $4=1+3$.

Критерий жёстче простого счёта размерностей. Кандидат обязан одновременно:
\begin{enumerate}
 \item содержать как минимум две ортогональные моды одного энергетического
 разрыва;
 \item быть независимым от охлаждаемого поперечного дублета;
 \item допускать физически разрешённый оператор связи;
 \item сохранять калибровочную, Real- и градуированную структуры;
 \item не повторять ранее закрытый двойной счёт.
\end{enumerate}

\section{Неожиданно положительная аффинная кратность три}

Аффинный KO6-проверка Тома V построил каноническое разложение
\begin{equation}
 \mathbb C^4=\operatorname{im}P_1\oplus\operatorname{im}P_3,
 \qquad
 P_1=\frac14J_4,
 \qquad
 P_3=I_4-P_1,
 \label{eq:v6-affine-sink-one-plus-three}
\end{equation}
с рангами один и три. На первом узле высота равна
\begin{equation}
 h_{\rm aff}=-P_3.
 \label{eq:v6-affine-sink-height}
\end{equation}
Следовательно, три неравномерные моды имеют в точности одну высоту и
образуют настоящее трёхмерное вырожденное подпространство.

Каноническая коизометрия
\begin{equation}
 V:\mathbb C^4\longrightarrow\mathbb C^3
\end{equation}
удовлетворяет
\begin{equation}
 VV^*=I_3,
 \qquad
 V^*V=P_3.
 \label{eq:v6-affine-sink-coisometry}
\end{equation}
Поэтому $V$ унитарно отождествляет аффинный триплет $P_3\mathbb C^4$ с
семейным триплетом. Требуемая резонансная кратность два не просто имеется:
она содержится в канонической кратности три и не требует $H_{15}$,
KO6 или нового поля.

\begin{proposition}[Резонансная кратность уже существует]
Аффинный первый узел содержит три ортогональные моды одной высоты и
каноническую полноранговую связь с семейным триплетом. Поэтому прежний
невырожденный потолок $2/3$ не применяется к полной аффинной архитектуре.
\end{proposition}

\section{Но это соседний угол, а не тензорный резервуар}

Полный частичный семейный модуль имеет структуру
\begin{equation}
 \mathbb C^4\oplus\mathbb C^3\oplus\mathbb C^3,
\end{equation}
а не $\mathbb C^4\otimes\mathbb C^3$. Поэтому аффинный триплет является
соседним узлом прямой суммы, связанным стрелкой $V$, но не независимой
подсистемой, для которой автоматически определён частичный след.

Можно, однако, рассматривать нормированное состояние физического угла
условно на том, что амплитуда осталась в среднем семейном узле. Пусть
выделенная ось остаётся на среднем узле, а две поперечные компоненты
частично переходят в $P_3$-угол. Если вероятность выживания каждой
поперечной компоненты равна $c^2$, условное состояние среднего угла имеет
спектр
\begin{equation}
 \frac{(1,c^2,c^2)}{1+2c^2}.
\end{equation}
Для точного состояния сосуществования требуется
\begin{equation}
 c^2=\frac{1-a_*}{2a_*}
 =0.0481454836\ldots,
 \label{eq:v6-affine-selective-survival}
\end{equation}
то есть угол поворота
\begin{equation}
 \theta_*=1.3495755726\ldots.
\end{equation}
При этом в семейном углу остаётся только
\begin{equation}
 q_G=\frac1{3a_*}=0.3654303224\ldots
\end{equation}
полной вероятности, а
\begin{equation}
 1-q_G=0.6345696776\ldots
\end{equation}
переходит в аффинный угол.

Таким образом, нужный условный спектр достижим энергосохраняющим
перераспределением между углами. Но это не прежнее охлаждение тензорной
подсистемы: вместе с упорядочением меняется вес самого физического угла.

\section{Каноническая связь остаётся изотропной}

Родительская стрелка имеет вид
\begin{equation}
 X=\rho V.
 \label{eq:v6-affine-canonical-isotropic-link}
\end{equation}
После отождествления $P_3\mathbb C^4\simeq\mathbb C^3$ она пропорциональна
тождеству семейного триплета. Поэтому одинаковый перенос трёх компонент
даёт
\begin{equation}
 R_G^{\rm cond}=I_3/3
\end{equation}
для любой общей амплитуды: изотропия сохраняется точно.

Чтобы получить \eqref{eq:v6-affine-selective-survival}, связь должна
различать ось и поперечную плоскость, например содержать
\begin{equation}
 X_P\sim V(I-P).
 \label{eq:v6-affine-state-selective-link}
\end{equation}
Но здесь $P$ уже является тем самым проектором, который механизм обязан
породить. Вставка $P$ в линейную связь была бы круговым выбором фазы.

\begin{nogocriterion}[Запрет заранее выбранной оси в стоке]
Нельзя считать аффинную кратность механизмом рождения $\mathbb{RP}^2$,
если оператор переноса заранее содержит проектор $P$. Допустима только
нелинейная или флуктуационно индуцированная обратная связь, для которой
изотропная точка сначала теряет устойчивость, а ось возникает как решение.
\end{nogocriterion}

\section{Почему остальные большие числа не являются стоком}

\subsection{Наблюдаемый пакет пятнадцать}

Разложение
\begin{equation}
 H_{15}=Q_L^{(6)}\oplus L_L^{(2)}\oplus u_R^{(3)}
 \oplus d_R^{(3)}\oplus e_R^{(1)}
\end{equation}
содержит пять попарно различных калибровочных представлений, каждое с кратностью
один. Размерности $6,2,3,3,1$ являются размерностями неприводимых зарядовых
пространств, а не числом эквивалентных копий. По лемме Шура калибровочные
инвариантный коммутант скалярен на каждом блоке и не может использовать
цвет или слабый изоспин как свободную память.

Это различие между размерностью представления и пространством кратности
стандартно в теории noiseless subsystems; см.
\path{doi:10.1103/PhysRevLett.79.3306} и
\path{doi:10.1103/PhysRevLett.84.2525}. Информация, доступная симметрично,
живет в факторе кратности, а не внутри единственной неприводимой копии.

\subsection{KO6 и ранг шесть модульного родителя}

KO6 добавляет сопряжённую половину, но она имеет противоположную
градуировку и Real-метку. Смешивание половин как двух взаимозаменяемых
каналов часов нарушает сохранение градуировки. Аналогично ранг шесть
вершинных проекторов $M_{18}$ раскладывается как семейный ранг три и его
$J$-сопряжение, а не как две свободные одинаковые копии внутри одного
физического сектора.

\subsection{Размерности матричных углов}

$M_{20}$ и $M_{35}$ являются матричными углами и связывающей алгеброй.
Их размер обозначает алгебру операторов, а не кратность собственных
состояний одного гамильтониана. Считать матричные единицы независимыми
принимающими частицами означало бы вновь смешать пространство операторов
с гильбертовым носителем.

\subsection{Real-пара внешнего квадрата}

Две одноориентационный копии уже были сведены физическим полуслед к одному
нормированному каналу. Повторное использование обеих как независимого
стока вернуло бы закрытый обменный двойной счёт.

\section{Вердикт: правильная кратность, неправильная обратная связь}

Аудит не дал простого ``нет''. Нужная энергетическая кратность уже
существует в самом раннем аффинном меню и канонически связана с семейным
триплетом. Это сильнее всех кандидатов $15$, $18$, $20$, $35$ и KO6.

Но линейная связь $\rho V$ сохраняет изотропию. Единственный минимальный
маршрут --- вывести нелинейную обратную связь, зависящую от состояния из уже
существующего функционала, не вставляя $P$ вручную. Следующая проверка должна
расширить амплитуду связи до наиболее общего $SO(3)$-эквивариантного
функционала от состояния $R$ и проверить, получает ли изотропная точка
отрицательную моду при сохранении общей энергии.

\begin{protocol}
\begin{enumerate}
 \item аффинное разложение $4=1+3$: \textbf{каноническое};
 \item вырожденная высота на $P_3$: \textbf{кратность три};
 \item связь $V:P_3\mathbb C^4\to\mathbb C^3$:
 \textbf{полноранговая};
 \item достаточная резонансная кратность: \textbf{найдена};
 \item независимый тензорный резервуар: \textbf{нет, соседний угол};
 \item условное достижение спектра $R_*$: \textbf{да};
 \item утечка из семейного угла: \textbf{$0.6345696776\ldots$};
 \item каноническая связь $\rho V$ создаёт расщепление $1+2$:
 \textbf{нет};
 \item $H_{15}$ как калибровочно-инвариантная кратность: \textbf{нет};
 \item KO6 как две свободные моды стока: \textbf{нет};
 \item $M_{20}/M_{35}$ как энергетическая кратность: \textbf{нет};
 \item следующая проверка:
 \textbf{нелинейная аффинная обратная связь};
 \item рождение материи:
 \textbf{резонансно переоткрыто, динамически не замкнуто}.
\end{enumerate}
\end{protocol}

Машинный сертификат сохранён в
\path{s2t_v6_existing_multiplicity_resonant_sink_gate_results.json}.